\SourcePage{224}{229}
\section{Radiation from diatomic molecules. Electronic spectra}

The distinctive features of molecular spectra arise primarily from the division of a molecule's energy into electronic, vibrational, and rotational parts, each successive part being small in comparison with the preceding one. The level structure of diatomic molecules was studied in detail in vol.~III, chap.~XI. Here we shall determine the resulting form of the spectrum and calculate the intensities of its lines\footnote{The discussion below is based on the material in vol.~III, \S\S~78, 82--88. To avoid burdening the text, we shall not refer continually to those sections.}.

We begin with the general case in which a transition changes the electronic state of the molecule (and, generally, its vibrational and rotational states as well). The frequencies of these transitions lie in the visible and ultraviolet regions of the spectrum. Their totality is called the molecule's \emph{electronic spectrum}. We shall always have electric-dipole transitions in mind; transitions of other types are generally of little importance in molecular spectroscopy.

As for dipole transitions in any system, the selection rule for the total angular momentum $J$ of the molecule is
\begin{equation}
|J'-J|\leq1\leq J+J'.
\tag{53.1}
\end{equation}
The strict selection rule for the parity of the system here corresponds to a selection rule for the sign of the level. Recall that, in the terminology adopted in molecular spectroscopy, states whose wavefunctions retain or reverse their sign under inversion (reversal of the coordinates of the electrons and nuclei) are called positive or negative, respectively. Thus the strict rule is
\begin{equation}
+\to-,\qquad -\to+.
\tag{53.2}
\end{equation}

If the molecule consists of identical atoms (with nuclei of the same isotope), its levels are also classified with respect to interchange of the nuclear coordinates: symmetric ($s$) levels, whose wavefunctions retain their sign under this transformation, and antisymmetric ($a$) levels, whose functions reverse sign. Since this transformation does not affect the electronic dipole-moment operator at all, its matrix elements are nonzero only for transitions that preserve this symmetry\footnote{This rule plainly applies also to transitions of any multipolarity.}:
\begin{equation}
s\to s,\qquad a\to a.
\tag{53.3}
\end{equation}

\SourcePage{225}{230}
This rule is not absolutely strict, however. The possession by a level of a given symmetry property is connected with the molecule having one particular value of the total nuclear spin $I$. Because the interaction of the nuclear spins with the electrons is extremely weak, $I$ is conserved to high accuracy, but not strictly. When this interaction is included, $I$ has no definite value, the symmetry property ($s$ or $a$) is not conserved, and the selection rule (53.3) is lost.

The electronic terms of a molecule made of identical atoms are also characterized by their parity ($g$ or $u$), that is, by the behavior of the wavefunctions when the signs of the electron coordinates (measured from the molecular center) are reversed while the nuclear coordinates remain fixed. This property of the electronic term is closely connected with the nuclear symmetry and the sign of the rotational levels belonging to that term. Levels belonging to an even ($g$) electronic term may have the characteristics $s+$ or $a-$, whereas those belonging to an odd term may have $s-$ or $a+$. Rules (53.2) or (53.3) therefore also imply
\begin{equation}
g\to u,\qquad u\to g.
\tag{53.4}
\end{equation}

As an approximate rule, (53.4) remains valid for molecules made of different isotopes of the same element. Since the nuclear charges are equal, if we consider the electronic term with fixed nuclei, we have a system of electrons in an electric field possessing a center of symmetry (at the point bisecting the internuclear distance). The symmetry of the electronic wavefunction under inversion through this point determines the parity of the term. Since the electric dipole-moment vector reverses sign under this transformation, rule (53.4) follows. A rule derived in this way is approximate because the nuclei have had to be treated as fixed. It is therefore violated when the interaction between the electronic state and molecular rotation is taken into account.

The subsequent selection rules depend on particular assumptions about the relative strengths of the various interactions in the molecule (that is, on its coupling case). These rules can consequently be only approximate.

Most electronic terms of diatomic molecules belong to coupling cases $a$ or $b$. Both are characterized by coupling of the orbital angular momentum to the axis (the electric interaction of the two atoms in the molecule) that is strong in comparison with all other interactions. There consequently exist quantum numbers $\Lambda$ and $S$ (the projection of the elec\SourcePageJoin{226}{231}tronic orbital angular momentum on the molecular axis and the total electron spin). The operator of the orbital quantity, the electronic dipole moment, commutes with the spin operator, so that
\begin{equation}
S'-S=0\qquad(\text{cases }a,b).
\tag{53.5}
\end{equation}
The change in $\Lambda$ obeys the selection rule
\begin{equation}
\Lambda'-\Lambda=0,\pm1\qquad(\text{cases }a,b),
\tag{53.6}
\end{equation}
and transitions between states with $\Lambda=0$ ($\Sigma$ terms) obey the additional rule
\begin{equation}
\Sigma^+\to\Sigma^+,\qquad \Sigma^-\to\Sigma^-\qquad(\text{cases }a,b)
\tag{53.7}
\end{equation}
(recall that the states $\Sigma^+$ and $\Sigma^-$ differ in their behavior under reflection in a plane through the molecular axis). Rules (53.6) and (53.7) are obtained by considering the molecule in a coordinate system rigidly attached to the nuclei (see III, \S~87); rule (53.6) is analogous to the selection rule for the magnetic quantum number in atoms.

Coupling cases $a$ and $b$ differ in the relation between the energy of the spin--axis interaction and the rotational energy (the separations of the rotational levels). The former energy is greater than the latter in case $a$, but much smaller in case $b$. We now consider the two cases separately.

\textbf{Case $a$.} Here there is a quantum number $\Sigma$, the projection of the total spin on the molecular axis (and with it the number $\Omega=\Sigma+\Lambda$, the projection of the total angular momentum). If both the initial and final states belong to case $a$, then
\begin{equation}
\Sigma'-\Sigma=0\qquad(\text{case }a)
\tag{53.8}
\end{equation}
(as follows from the already mentioned commutation of the dipole moment with the spin). Equations (53.6) and (53.8) imply\footnote{This rule also remains valid in case $c$ (where coupling of the orbital angular momentum to the axis is weak in comparison with spin--orbit coupling), in which the numbers $\Lambda$ and $\Sigma$ do not exist separately.}
\begin{equation}
\Omega'-\Omega=0,\pm1.
\tag{53.9}
\end{equation}
If $\Omega=\Omega'=0$, transitions with $J'=J$ are forbidden in addition to the general rule (53.1)\footnote{This rule is analogous to the prohibition of transitions with $J=J'$ when $M=M'=0$ in atoms (see \S~5), where, however, it could be relevant only in the presence of an external field. Here the rule follows directly from formula (53.12) below; the $3j$ symbol $\begin{pmatrix}J&1&J\\0&0&0\end{pmatrix}$ vanishes because the sum $J+J+1$ is odd.}:
\begin{equation}
J'-J=\pm1\quad\text{for }\Omega=\Omega'=0\qquad(\text{case }a).
\tag{53.10}
\end{equation}

\SourcePage{227}{232}
Consider transitions between any two specified vibrational levels belonging to two different electronic terms (of case $a$). When the fine structure of an electronic term is taken into account, each of these levels splits into several components; by rule (53.5), their number ($2S+1$) must be the same for both levels. According to rule (53.8), each component of one level combines with only one component of the other, having the same value of $\Sigma$.

Now take one pair of levels with the same $\Sigma$; their values of $\Omega$ and $\Omega'$ may differ (together with $\Lambda$ and $\Lambda'$) by 0 or $\pm1$. When rotation is included, each splits into a sequence of levels distinguished by the numbers $J$ and $J'$, which take the values $J\geq|\Omega|$ and $J'\geq|\Omega'|$. The dependence of the transition probabilities on these numbers can be found in general form (H.~Hönl, R.~London, 1925).

The matrix element for the transition $n\Lambda\Omega JM_J\to n'\Lambda'\Omega'J'M'_J$ (where $n$ denotes the remaining characteristics of the electronic term, apart from $\Omega$ and $\Lambda$) is
\begin{equation}
\begin{aligned}
&|\langle n'\Lambda'\Omega'J'M'_J|d_q|n\Lambda\Omega JM_J\rangle|={}\\
&\quad=\sqrt{(2J+1)(2J'+1)}
\begin{pmatrix}J'&1&J\\-\Omega'&q'&\Omega\end{pmatrix}
\begin{pmatrix}J'&1&J\\-M'_J&q&M_J\end{pmatrix}
|\langle n'\Lambda'|\bar d_{q'}|n\Lambda\rangle|,
\end{aligned}
\tag{53.11}
\end{equation}
where $d_q$ and $\bar d_{q'}$ are the spherical components of the dipole-moment vector in the fixed coordinate system $xyz$ and in the ``moving'' system $\xi\eta\zeta$, respectively, with the $\zeta$ axis along the molecular axis (this formula follows with the aid of formula (110.6), see III). The matrix elements $\langle n'\Lambda'|\bar d_{q'}|n\Lambda\rangle$ are independent of the rotational quantum numbers $J,J'$ and depend only on the characteristics of the electronic terms (in the present case they are also independent of $\Sigma$\footnote{This may be verified in the same way as at the beginning of III, \S~29 for a scalar quantity $f$. Here the vector operator $\mathbf d$ commutes with the operator of the vector $\mathbf S$, which is conserved in the zeroth approximation, while $\Sigma$ is the projection of $\mathbf S$ on the $\zeta$ axis in the rotating coordinate system in which the commutation condition for $\mathbf d$ and $\mathbf S$ must be considered.}); the indices $\Omega'=\Lambda'+\Sigma$ and $\Omega=\Lambda+\Sigma$ have therefore been omitted from the notation for the matrix element.

The probability of the transition $n\Lambda\Omega J\to n'\Lambda'\Omega'J'$ is proportional to the square of the matrix element (53.11), summed over $M'_J$. By formula (106.12) (see III),
\[
\sum_{M'_J}\begin{pmatrix}J'&1&J\\-M'_J&q&M_J\end{pmatrix}^{\!2}=\frac1{2J+1},
\]
\SourcePage{228}{233}
and hence
\begin{equation}
\begin{aligned}
w(n\Lambda\Omega J\to n'\Lambda'\Omega'J')={}&(2J'+1)
\begin{pmatrix}J'&1&J\\-\Omega'&\Omega'-\Omega&\Omega\end{pmatrix}^{\!2}\\
&\times B(n',n;\Lambda',\Lambda),
\end{aligned}
\tag{53.12}
\end{equation}
where the coefficients $B$ are independent of $J,J'$ (we of course neglect the relatively small difference among the frequencies of transitions with different $J,J'$)\footnote{When $\Lambda$ doubling is included, each rotational level $J$ considered here splits further into two levels, one positive and the other negative. Thus, in place of one transition $J\to J'$, rule (53.2) gives two transitions: from the positive (negative) component of level $J$ to the negative (positive) component of level $J'$. The probabilities of these transitions are equal.}.

Summing (53.12) over $J'$ gives simply $B(n',n;\Lambda',\Lambda)$, by the orthogonality property of the $3j$ symbols (see III, (106.13)). In other words, the total probability of transition from a rotational level $J$ of the state $\Omega$ to all levels $J'$ of the state $\Omega'$ is independent of $J$.

\textbf{Case $b$.} In this case, besides the total angular momentum $J$, there is also a quantum number $K$, the angular momentum of the molecule excluding its spin. The selection rules for this number coincide with the general selection rules for any orbital vector quantity (such as the electric dipole moment):
\begin{equation}
|K'-K|\leq1\leq K+K'\qquad(\text{case }b)
\tag{53.13}
\end{equation}
with the additional prohibition of the transition $K=K'$ when $\Lambda=\Lambda'=0$ (analogously to (53.10)):
\begin{equation}
K'-K=\pm1\quad\text{for }\Lambda=\Lambda'=0.
\tag{53.14}
\end{equation}

Consider transitions between rotational components of specified vibrational levels of two electronic states belonging to case $b$. Their transition probabilities are given by the same formulas (53.12), with $K,\Lambda$ written in place of $J,\Omega$. When the fine structure is included (for $S\ne0$), each rotational level $K$ splits into $2S+1$ components: $J=|K-S|,\ldots,K+S$, so that a multiplet appears in place of a single line $K\to K'$. Since this case involves the addition of the free angular momenta $K$ and $S$ (which are not coupled to the molecular axis), the formulas for the relative probabilities of the various multiplet lines coincide with the analogous formulas (49.15) for fine-structure components of atomic spectra, where the angular momenta $L$ and $S$ play the same role (in the case of $LS$ coupling).

\SourcePage{229}{234}
We have thus considered the selection rules that determine the possible spectral lines in all the principal cases encountered in diatomic molecules.

The set of lines from transitions between the rotational components of two given electronic-vibrational levels constitutes what spectroscopy calls a \emph{band}: because the rotational intervals are small, the lines in a band lie very close together. Their frequencies are given by the differences
\begin{equation}
\hbar\omega_{J,J'}=\mathrm{const}+BJ(J+1)-B'J'(J'+1),
\tag{53.15}
\end{equation}
where $B,B'$ are the rotational constants in the two electronic states (to avoid needless complications, we assume singlet electronic terms). For $J'=J,J\pm1$, formula (53.15) is represented graphically (Fig.~2) by three branches (parabolas), whose points at integer $J$ give the frequencies (the arrangement of the branches in Fig.~2 corresponds to $B'<B$; for $B'>B$ they open toward smaller $\omega$, and the curve for $J'=J-1$ lies uppermost)\footnote{The series of lines corresponding to transitions with $J'=J+1,J,J-1$ are called the $P$, $Q$, and $R$ branches, respectively.}. As is clear from the figure, the presence of a turning branch causes the lines to crowd toward a particular limiting position (the band head).

\begin{figure}[ht]
\centering
\begin{tikzpicture}[x=.58cm,y=.42cm,>=Stealth,line width=.55pt,font=\small]
  \draw[->] (0,0) -- (11.2,0) node[right] {$\omega$};
  \draw[->] (0,0) -- (0,10.5) node[above] {$J$};
  \foreach \j in {0,...,9} {
    \draw (-.12,\j) -- (.12,\j);
    \ifodd\j\else\node[left] at (-.18,\j) {\j};\fi
  }
  \draw[smooth,domain=0:9,samples=60] plot ({2.15+.085*\x*\x},{\x});
  \draw[smooth,domain=0:9,samples=60] plot ({4.45+.075*(\x-1.8)*(\x-1.8)},{\x});
  \draw[smooth,domain=1:9,samples=60] plot ({8.1-.33*\x+.058*\x*\x},{\x});
  \foreach \j in {1,...,9} {
    \draw[densely dashed,gray] (0,\j) -- ({8.1-.33*\j+.058*\j*\j},\j);
  }
  \draw[densely dashed,gray] (4.45,0) -- (4.45,10);
  \node[anchor=west] at (8.1,9.2) {$J'=J+1$};
  \node[anchor=west] at (7.5,6.6) {$J'=J$};
  \node[anchor=west] at (5.9,3.4) {$J'=J-1$};
\end{tikzpicture}
\caption{}
\end{figure}

In discussing line intensities, we must also mention the distinctive alternation of intensities in certain bands of the electronic spectra of molecules made of atoms of the same isotope (W.~Heisenberg, F.~Hund, 1927). The symmetry requirements associated with the nuclear spins cause rotational components with even and odd values of $K$ in electronic $\Sigma$ terms to have opposite nuclear symmetry and, accordingly, different nuclear statistical weights $g_s$ and $g_a$ (see III, \S~86). By rule (53.14), only $J'=J\pm1$ is allowed for transitions between two different $\Sigma$ terms; by (53.4), one of the $\Sigma$ terms must be even and the other odd. Consequently, for a specified value of $J'-J$, transitions at successive values of $J$ occur alternately between pairs of symmetric levels and pairs of antisymmetric levels (as illustrated in Fig.~3 for the states $\Sigma_u^+$ and $\Sigma_g^+$). On the other hand, the observed line intensity is propor\SourcePageJoin{230}{235}tional to the number of molecules in the given initial state and therefore to its statistical weight. Hence the intensities of successive lines ($J=0,1,2,\ldots$) will alternate between larger and smaller values, proportional in turn to $g_s$ and $g_a$ (apart from the monotonic variation predicted by formulas (53.12))\footnote{Here it is assumed that all states with different values of the total nuclear spin are uniformly populated.}.

\begin{figure}[ht]
\centering
\begin{tikzpicture}[x=1.15cm,y=.9cm,>=Stealth,line width=.55pt,font=\small]
  \node[left] at (-.6,2.1) {$\Sigma_u^+$};
  \node[left] at (-.6,0) {$\Sigma_g^+$};
  \foreach \j in {0,...,5} {
    \node at (\j,2.8) {$\j$};
    \node at (\j,-.7) {$\j$};
  }
  \node[left] at (-.55,2.8) {$J=$};
  \node[left] at (-.55,-.7) {$J'=$};
  \foreach \j/\sym/\sgn in {0/a/+,1/s/-,2/a/+,3/s/-,4/a/+,5/s/-} {
    \node at (\j,2.48) {$\sym$};
    \node[draw,circle,fill=white,inner sep=0pt,minimum size=3.1mm] at (\j,2.1) {$\scriptstyle\sgn$};
  }
  \foreach \j/\sym/\sgn in {0/s/+,1/a/-,2/s/+,3/a/-,4/s/+,5/a/-} {
    \node[draw,circle,fill=white,inner sep=0pt,minimum size=3.1mm] at (\j,0) {$\scriptstyle\sgn$};
    \node at (\j,-.35) {$\sym$};
  }
  \foreach \j in {0,...,4} \draw (\j,1.94) -- (\j+1,.16);
  \foreach \j in {1,...,5} \draw[densely dashed] (\j,1.94) -- (\j-1,.16);
\end{tikzpicture}
\caption{}
\end{figure}

There are no strict selection rules for the change of the vibrational quantum number in transitions between two different electronic terms. There is, however, a rule (the Franck--Condon principle) that predicts the most probable change of vibrational state. It rests on the quasiclassical character of the nuclear motion resulting from the large nuclear masses (compare the discussion of predissociation in III, \S~90)\footnote{Strictly speaking, the vibrational quantum number must also be sufficiently large.}.

\begin{figure}[ht]
\centering
\begin{tikzpicture}[x=.9cm,y=.72cm,>=Stealth,line width=.6pt,font=\small]
  \draw[->] (.55,.25) -- (.55,5.95);
  \node[left] at (.48,5.72) {$U(r)$};
  \node[left] at (.48,3.15) {$U'(r)$};
  \draw[smooth,domain=.56:4.05,samples=70] plot ({\x},{3.75+.55*(\x-1.75)*(\x-1.75)});
  \draw[smooth,domain=.56:4.05,samples=70] plot ({\x},{1.05+.36*(\x-1.65)*(\x-1.65)});
  \draw (.72,4.18) -- (2.72,4.18) node[right] {$E$};
  \draw (.72,1.62) -- (2.72,1.62);
  \node[above left] at (2.55,1.62) {$E'$};
  \draw[densely dashed] (2.72,1.62) -- (2.72,4.18);
  \draw[->] (2.94,4.12) -- (2.94,1.68);
\end{tikzpicture}
\caption{}
\end{figure}

In the integral determining the matrix element of the transition between the vibrational states $E$ and $E'$ of the electronic terms $U(r)$ and $U'(r)$, the main contribution comes from the neighborhood of the point $r=r_0$ at which
\begin{equation}
U(r_0)-U'(r_0)=E-E'
\tag{53.16}
\end{equation}
(that is, the momenta of the relative nuclear motion are the same in both states: $p=p'$). For a specified $E$, the transition probability (as a function of the final energy $E'$) is the greater, the smaller each of the differences $E-U$ and $E'-U'$. It is greatest when
\begin{equation}
E-U(r_0)=E'-U'(r_0)=0,
\tag{53.17}
\end{equation}
that is, when the ``transition point'' $r_0$ (a root of equation (53.16)) coincides with a classical turning point of the nuclei (Fig.~4 illustrates graphically this relation between $E$ and the most probable $E'$). In pictorial terms, the transition is most probable \SourcePage{231}{236}near a point where the nuclei stop and in whose neighborhood they consequently spend more time.
